\section{Experimental design}
\label{sec:design}

\subsection{Predictions}

Each prediction below was chosen against one filter: \emph{would a careful reader bet
on it in advance?} Claims that pass that filter trivially---an open network produces
more variety, a persistent team completes more---are excluded, because confirming
them consumes a study and returns nothing.

\begin{enumerate}[label=\textbf{H\arabic*},ref=\arabic*,leftmargin=2.2em,itemsep=2pt]
\item\label{h:arch} \textbf{Architecture dominates text.}
      $|\Delta_{\text{architectural}}| \gg |\Delta_{\text{textual}}|$ on the primary
      metrics. \emph{Falsified if} seeded trust, shared goals, accountability or
      tenure move outcomes comparably to the architectural knobs.

\item\label{h:reach} \textbf{Equalizing reach erases most of the gap.}
      Comparing $A$ to $D$ overstates the public/private difference relative to the
      isolated $\Delta_{\text{access}}$. \emph{Falsified if} the two main effects are
      of similar magnitude.

\item\label{h:lossy} \textbf{Extraction is lossy; delegation is not.}
      On decomposable subtasks a sandboxed counterpart beats a readable one, because
      answering uses the holder's full local context whereas reading forces the
      requester to re-join partial material. \emph{Falsified if} store access
      dominates uniformly.

\item\label{h:afford} \textbf{Agents do not use the affordances they are given.}
      With \texttt{list} and \texttt{read} available, requesters continue to ask
      questions; with persistent memory available, they do not reuse it. \emph{
      Falsified if} tool-use distributions shift materially when the affordance
      appears.

\item\label{h:pollute} \textbf{The trusted channel absorbs more misinformation.}
      Distractor absorption is \emph{higher} under store access than under sandboxed
      access, because a question forces the holder to filter while a read does not.
      \emph{Falsified if} absorption tracks distractor exposure alone.

\item\label{h:inert} \textbf{Optima do not move under seeding.}
      $N^{*}$ is unchanged between seeded-trust and control. \emph{Falsified if}
      $N^{*}$ shrinks under seeded trust---which would mean the seed changed the
      requester's stopping rule, and hence that relational text is not inert.

\item\label{h:cross} \textbf{Crossing has an interior optimum.}
      Hybrid utility is non-monotone in \metric{crossing\_count}: zero crossings
      under-explore, many crossings lose context and pay reconciliation repeatedly.
      \emph{Falsified if} utility is monotone in either direction.
\end{enumerate}

\subsection{Conditions}

Three task conditions, each realized as configurations of the same kernel: public
strangers ($A$), a pre-existing private team ($D$), and a private team with public
consultation. The third is not recruitment---public counterparts never gain
membership, presence, or execution authority. Material returns only by being
distilled into a file under the private tree with provenance, so that reading a
public answer and writing a private file remain separate permission decisions. Cells
$B$ and $C$ complete the factorial and are what allow the two main effects to be
separated.

\subsection{Scale variables as probes}

We report $N$ (public query breadth), $K$ (private team size), $T$ (interaction
history), $Q$ (consultation frequency) and $R$ (fraction of external material
imported). We do not report an optimum as a finding on its own: an $N^{*}$ obtained in
one setup is setup-specific and generalizes poorly.

Instead each is measured \emph{under both seeded and control conditions}, and the
quantity of interest is the difference between the two curves. Under \hyp{h:inert} an
optimum is a readout of whether a relational property changed a policy: if seeded
trust does not shift where a requester stops asking, the seed did not enter the
stopping rule. $Q$ and $R$ receive the most attention because they concern the
boundary, which is the paper's subject; $N$ and $K$ are reported as probes and
supporting curves.

\subsection{Controls and preregistration}

One model across all agents and all cells; one runtime; one kernel. Identical task
prompts, deliverable schemas, scoring functions, token budgets, and seed
distributions. Rate limits and loop bounds are logged per episode rather than absorbed
into results. Rubrics, schemas, normalization rules and the scorer signature are
frozen before any run; the mapping from each hypothesis to the figure that tests it is
fixed in advance and reproduced in the appendix.

\paragraph{Preliminary observation.}
A pilot conducted before the access axis was implemented compared open against
bounded discovery with both configurations exposing readable stores---that is, cells
$B$ and $D$ rather than $A$ and $D$. Equalizing reach in that pilot removed the large
majority of the measured difference. We report it in
Section~\ref{sec:results} as a bound on $\Delta_{\text{discovery}}$ and as the
motivation for the factorial, and we note its limitation: with both cells on the
readable column it could not have detected an access effect at all.
